% Generated draft: requires mathematical and source review.
\documentclass[11pt]{article}
\usepackage[margin=25mm]{geometry}
\usepackage{fontspec}
\setmainfont{Latin Modern Roman}
\usepackage{amsmath,amssymb,mathtools,unicode-math}
\setmathfont{Latin Modern Math}
\usepackage{enumitem,needspace,xurl,hyperref}
\setlength{\parindent}{0pt}
\setlength{\parskip}{5pt}
\setlength{\emergencystretch}{3em}
\newcommand{\N}{\mathbb{N}}
\newcommand{\Z}{\mathbb{Z}}
\newcommand{\Q}{\mathbb{Q}}
\newcommand{\R}{\mathbb{R}}
\newcommand{\C}{\mathbb{C}}
\newcommand{\F}{\mathbb{F}}
\newcommand{\A}{\mathbb{A}}
\newcommand{\PP}{\mathbb{P}}
\newcommand{\E}{\mathbb{E}}
\newcommand{\eps}{\varepsilon}
\newcommand{\dd}{\,\mathrm{d}}
\begin{document}
\section*{UnsolvedMath ID: 30005775}
\textbf{Source problem code:} catalogue entry
\subsection{Definitions and mathematical statement}
Work over the field \(\C\). Let \(G\) be a group, with no finiteness assumption in the supplied statement. A left-left Yetter--Drinfeld module \(V\in{}_{\C G}^{\C G}\mathcal{YD}\) is a \(G\)-graded complex vector space \(V=\bigoplus_{g\in G}V_g\) with a left \(G\)-action satisfying
\[
h\cdot V_g\subseteq V_{hgh^{-1}}\qquad(g,h\in G).
\]
Its braiding is \(c(v\otimes w)=g\cdot w\otimes v\) for \(v\in V_g\). The Nichols algebra \(\mathcal B(V)\) is the quotient of the braided tensor Hopf algebra \(T(V)\), whose elements of \(V\) are primitive, by the largest homogeneous braided Hopf ideal contained in \(\bigoplus_{n\ge2}V^{\otimes n}\).

For this draft, a post-Nichols algebra of \(V\) means a connected nonnegatively graded braided Hopf algebra
\[
R=\bigoplus_{n\ge0}R_n\quad\text{in }{}_{\C G}^{\C G}\mathcal{YD},\qquad R_0=\C 1,
\]
with a specified identification \(R_1\cong V\) and with
\[
P(R):=\{x\in R:\Delta(x)=x\otimes1+1\otimes x\}=R_1.
\]
Under this convention, the braided Hopf subalgebra generated by \(R_1\) is the canonical copy of \(\mathcal B(V)\). The supplied excerpt does not explicitly define post-Nichols algebra, so this convention requires comparison with the full report.

\textbf{Question.} For every group \(G\) and every \(V\in{}_{\C G}^{\C G}\mathcal{YD}\) such that \(\dim_{\C}\mathcal B(V)<\infty\), is every finite-dimensional post-Nichols algebra \(R\) of \(V\) generated as an algebra by \(R_1\)? Equivalently, must the canonical inclusion
\[
\mathcal B(V)\hookrightarrow R
\]
be an isomorphism of graded braided Hopf algebras respecting the specified degree-one identification?

\textbf{Scope qualification.} The supplied mathematical statement allows arbitrary groups, whereas the supplied literature assessment discusses arbitrary finite groups and the surrounding excerpt concerns finite-dimensional pointed Hopf algebras. This draft retains the stated arbitrary-group quantifier; it does not silently restrict \(G\) to be finite.

\textbf{Source context.} The excerpt places this question, described there as generation in degree one, among subproblems of a conjecture about finite-dimensional pointed Hopf algebras. For such an algebra \(H\), with group of group-like elements \(G(H)\) and infinitesimal braiding \(V\), that surrounding conjecture asserts
\[
H\simeq\bigl(\mathcal B(V)\#\C G(H)\bigr)_{\sigma},
\]
where \(\#\) denotes bosonization and \(\sigma\) is an invertible Hopf \(2\)-cocycle on the bosonization with values in \(\C\). This is supplied conjectural context, not an additional assertion or a replacement for the question above.
\subsection{Short English statement}
Over \(\C\), suppose the Nichols algebra of a Yetter--Drinfeld module \(V\) over a group \(G\) is finite-dimensional. Must every finite-dimensional post-Nichols algebra with degree-one part \(V\) equal its Nichols subalgebra, and hence be generated in degree one? The supplied statement allows arbitrary groups; the finite-group emphasis of the literature assessment and the precise post-Nichols convention require review.
\subsection{Sources}
\begin{itemize}
\item Local UnsolvedMath registry: supplied complete source record for permanent ID \(30005775\), titled \emph{Uniqueness of Finite-Dimensional Post-Nichols Algebras}, exact source code \texttt{catalogue entry}, in the Oberwolfach Reports Open Problems collection. Project registry locations are \texttt{data/unsolvedmath/problems.json} and \texttt{catalog/problems.json}; these files were not independently inspected for this draft.
\item Supplied primary citation: \emph{Mini-Workshop: Bridging Number Theory and Nichols Algebras via Deformations} (2024), Oberwolfach Reports, pp.\ 235--272, DOI \texttt{10.4171/owr/2024/5}. Source URL and literature-source link labelled ``original OWR report'': \url{https://doi.org/10.4171/owr/2024/5}. The record supplies no individual proposer.
\item Additional supplied literature source: \emph{Lifting via cocycle deformation}, \url{https://arxiv.org/abs/1212.5279}. No authors or more specific location supporting the assessment are supplied in the record.
\item Supplied literature assessment, checked \textbf{2026-08-22}: the record labels the problem open and reports that generation in degree one is known for abelian groups and documented cases, while the question for arbitrary finite groups with finite-dimensional Nichols algebras remains unresolved. This is supplied dated curation, explicitly described as literature triage derived from the OWR curation; neither its status nor its cited claims have been independently verified here.
\item The registry labels the statement \texttt{corrected\_verified} and reports that it was reconstructed and checked against the official Oberwolfach report. This is a supplied verification claim, not an independent source check performed for this draft.
\end{itemize}
\end{document}
